\documentclass[10pt,a4paper]{amsart}
\usepackage[margin=19mm]{geometry}
\usepackage{amsmath,amssymb,mathtools}
\usepackage[scr=rsfs,cal=euler]{mathalfa}
\usepackage{xfrac}
\usepackage[unicode,hidelinks,bookmarks=false]{hyperref}
\newcommand{\ie}{i.e.\ }
\newcommand{\cf}{cf.\ }
\newcommand{\resp}{respectively\ }
\newcommand{\Subsection}{\subsection}
\providecommand{\nobreakdash}{\nobreak\hbox{-}}
\setlength{\emergencystretch}{2em}
\multlinegap0pt
\usepackage{bm}
\usepackage{bbm}
\usepackage{dsfont}
\usepackage{xspace}
\usepackage{cancel}
\usepackage{mathtools}
\usepackage{amsthm}
\makeatletter
\@ifundefined{theorem}{\newtheorem{theorem}{Theorem}}{}
\@ifundefined{proposition}{\newtheorem{proposition}[theorem]{Proposition}}{}
\makeatother
\title[A Probabilistic Interpretation of the Ball Mapper Graph]{A Probabilistic Interpretation of the Ball Mapper Graph}
\author{John Rick Manzanares, Jay-Anne Bulauan}
\date{Source: arXiv:2608.11397v1 (CC BY 4.0)}
\begin{document}
\maketitle

\noindent\emph{Source and licence.} A Probabilistic Interpretation of the Ball Mapper Graph. Version arXiv:2608.11397v1 (CC BY 4.0), distributed under the Creative Commons Attribution 4.0 International licence, as recorded in the arXiv licence metadata for this exact version and shown on the abstract page https://arxiv.org/abs/2608.11397v1. Full rights notice: licenses/arxiv-2608.11397-CC-BY-4.0.txt.\par\smallskip

\noindent\textit{Editorial scope.} Counted unit: the Proposition whose source label is \texttt{None} in the article (source0.tex lines 342--344, source label \texttt{None}), delivered together with the article's own complete original proof of it (source0.tex lines 346--369, one \texttt{proof} environment of 24 lines). No mathematics is written, repaired or re-derived: statement and proof are the article's own text line for line. Required context retained uncounted: none. Every definition, display and label of those lines is retained unchanged. Omitted from this package and not redistributed: 1--341, 345--345, 370--833. Nothing inside a retained excerpt is omitted or altered: every retained body is delivered exactly as the article's lines are written, so no omission receipt of any kind accompanies this package. Citations: the retained excerpts carry no citation command at all, so no bibliography entry of the article is transcribed and no reference is redistributed. Reference adaptation: none was needed and none was made; every reference inside the delivered unit resolves inside it, so no unresolved reference or citation is produced. Printed slips kept literal: no printed word, symbol, hypothesis or number of the counted statement or of its proof was altered. Layout and class: the article's own class is \texttt{extarticle} with packages including babel, fontenc, stix2, FiraSans, amsmath, amsthm, geometry, xcolor, microtype, authblk, enumitem, fancyhdr, lastpage, lettrine; the wrapper is the lane's pinned \texttt{amsart} editorial wrapper at 10pt on A4, which restates the theorem-like environments the retained text needs and adds the article's own macro definitions that the retained text needs (none), with extra wrapper packages (bm, bbm, dsfont, xspace, cancel, mathtools). The delivered numbering is the unit's own. Assets: no retained span contains a figure environment, a graphics-inclusion command, a PDF-page inclusion, a table or a TikZ picture; no image file is copied into this package, the package declares no asset dependency and no image file is redistributed with it. Licence: version arXiv:2608.11397v1 of this article is distributed under the Creative Commons Attribution 4.0 International licence, as recorded in the arXiv licence metadata for this exact version and shown on the abstract page https://arxiv.org/abs/2608.11397v1. A full rights notice is delivered at licenses/arxiv-2608.11397-CC-BY-4.0.txt. Characters: every retained excerpt is pure ASCII, so the delivered engine, XeLaTeX with its default Latin Modern text font, reports no missing character.
\begin{proposition}
For every $x\in X$, the function $\kappa_x$ defined above is a probability measure on $\mathcal C_\varepsilon(X)$.
\end{proposition}

\begin{proof}
Fix $x\in X$ and write
\[
Z_x=\sum_{j=1}^K a_j(x).
\]
By assumption, $Z_x > 0$. Since each $a_i(x)$ is nonnegative, we have
\[
0\le \kappa_x(A)\le 1
\]
for every $A\subseteq\mathcal C_\varepsilon(X)$. Also, $\kappa_x(\emptyset)=0$ and
\[
\kappa_x(\mathcal C_\varepsilon(X)) = \frac{\sum_{i=1}^K a_i(x)}{Z_x} = 1.
\]

Now let $A_1,\dots,A_m\subseteq \mathcal C_\varepsilon(X)$ be pairwise disjoint. Then
\[
\sum_{\left\{i \ \mid \ C_i\in \bigcup_{\ell=1}^m A_\ell\right\}} a_i(x) =\sum_{\ell=1}^m \sum_{\{i \ \mid \ C_i\in A_\ell\}} a_i(x)
\]
because no ball belongs to more than one of the sets $A_\ell$. Dividing by $Z_x$ gives
\[
\kappa_x\left(\bigcup_{\ell=1}^m A_\ell\right) = \sum_{\ell=1}^m \kappa_x(A_\ell).
\]
Thus, $\kappa_x$ is a probability measure on $\mathcal C_\varepsilon(X)$.
\end{proof}

\end{document}
